\chapter{Configuration and Spatial Verification of the Baselines}
\label{app:baseline_configurations}

To ensure complete scientific reproducibility and to prevent the disruption of the narrative flow in Chapter~\ref{ch:fenics_implementation}, all comprehensive configuration matrices, domain truncation proofs, and spatial grid independence studies for the fluid benchmarks are documented in this appendix.

\section{Verification Benchmark 1: Channel Flow}
\label{sec:app_channel}

\subsection{Physical and Geometric Parameters}
\label{subsec:app_channel_physical}
The first verification benchmark consists of a pressure-driven periodic 2D channel. To align with fundamental turbulence research conventions, the characteristic length is defined as the full channel height ($H$), and the Reynolds number ($Re_H$) is computed strictly based on this geometric parameter rather than the hydraulic diameter. The complete physical setup is detailed in Table~\ref{tab:app_channel_physical}.

\begin{table}[htbp]
	\centering
	\footnotesize
	\caption{Geometric, physical, and boundary-condition parameters for the periodic channel flow benchmark.}
	\label{tab:app_channel_physical}
	\renewcommand{\arraystretch}{1.2}
	\begin{tabularx}{\textwidth}{l X}
		\hline
		\textbf{Parameter} & \textbf{Channel flow setting} \\
		\hline
		Benchmark type & Pressure-driven, streamwise-periodic 2D channel \\
		Domain dimensions & Length $L=2.0$ m; full channel height $H=2.0$ m \\
		Characteristic length & Full channel height $H=2.0$ m \\
		Reference velocity & $U_{\mathrm{ref}}=18.5$ m/s, utilized to define the target Reynolds number \\
		Kinematic viscosity & $\nu = 1.81818 \times 10^{-3}$ m$^2$/s \\
		Reynolds number & $Re_H \approx 2.03 \times 10^4$ \\
		Body force & $\mathbf{f}=(0,0)$ m/s$^2$ \\
		Driving condition & Fixed pressure drop from $p=2$ Pa to $p=0$ Pa across the periodic streamwise direction \\
		Velocity boundary conditions & Periodic velocity between inlet/outlet planes; no-slip on top and bottom walls \\
		Pressure boundary conditions & $p=2$ Pa at the upstream pressure plane; $p=0$ Pa at the downstream pressure plane \\
		SA boundary conditions & Periodic $\tilde{\nu}$ between inlet/outlet planes; $\tilde{\nu}=0$ at walls \\
		Initial fields & $\mathbf{u}_0=(18.5,0)$ m/s, $p_0=2$ Pa, $\tilde{\nu}_0=5.0 \times 10^{-3}$ m$^2$/s \\
		Verification mesh & \texttt{Fine\_WallResolved}; first layer $\Delta y_1 = 1.709 \times 10^{-3}$ m \\
		\hline
	\end{tabularx}
\end{table}

\subsection{Numerical Simulation Parameters}
\label{subsec:app_channel_numerical}
The channel flow was resolved using an Incremental Pressure Correction Scheme (IPCS) embedded within a Picard fixed-point iteration loop. The exact numerical solver configurations, convergence tolerances, and relaxation factors utilized in FEniCS are summarized in Table~\ref{tab:app_channel_numerical}.

\begin{table}[htbp]
	\centering
	\footnotesize
	\caption{Numerical simulation parameters for the FEniCS steady Spalart-Allmaras channel benchmark.}
	\label{tab:app_channel_numerical}
	\renewcommand{\arraystretch}{1.2}
	\begin{tabularx}{\textwidth}{l c X}
		\hline
		\textbf{Parameter} & \textbf{Symbol} & \textbf{Setting} \\
		\hline
		Velocity space & $\mathbf{u}$ & Vector CG2, periodic in streamwise direction \\
		Pressure space & $p$ & Scalar CG1 \\
		SA working-variable space & $\tilde{\nu}$ & Scalar CG1, periodic in streamwise direction \\
		Quadrature degree & -- & $2$ \\
		Maximum outer Picard steps & $N_{\mathrm{Picard,max}}$ & $200$ \\
		SA solves per Picard step & -- & $1$ \\
		Outer velocity tolerance & $\epsilon_{\mathbf{u}}$ & $1.0 \times 10^{-6}$ \\
		Outer pressure tolerance & $\epsilon_p$ & $1.0 \times 10^{-6}$ \\
		Outer SA tolerance & $\epsilon_{\tilde{\nu}}$ & $1.0 \times 10^{-6}$ \\
		IPCS pseudo-time step & $\Delta t$ & $5.0 \times 10^{-3}$ s \\
		Maximum IPCS steps per Picard step & -- & $500$ \\
		IPCS velocity tolerance & -- & $1.0 \times 10^{-6}$ \\
		IPCS pressure tolerance & -- & $1.0 \times 10^{-6}$ \\
		Velocity relaxation factor & $\alpha_u$ & $0.3$ \\
		Pressure relaxation factor & $\alpha_p$ & $0.3$ \\
		SA relaxation factor & $\alpha_{\tilde{\nu}}$ & $0.3$ \\
		SA lower bound & $\tilde{\nu}_{\min}$ & $1.0 \times 10^{-12}$ \\
		Wall-distance equation & -- & Relaxed Eikonal, $\sigma_w=0.1$, $G_0=20$ \\
		Wall-distance Newton settings & -- & Relative tolerance $1.0 \times 10^{-8}$; absolute tolerance $1.0 \times 10^{-10}$; maximum $80$ iterations; relaxation $0.5$ \\
		Linear solver strategy & -- & IPCS blocks: direct MUMPS; SA transport: default solver \\
		\hline
	\end{tabularx}
\end{table}

\subsection{Spatial Verification and Grid Independence}
\label{subsec:app_channel_mesh}
To guarantee that the numerical results were independent of the spatial discretization, a formal grid independence study was conducted. Three distinct boundary-layer inflated meshes (Coarse, Medium, and Fine) were evaluated. The geometric properties, node counts, and resolution parameters of this mesh hierarchy are explicitly detailed in Table~\ref{tab:appendix_channel_meshes}. Crucially, the first-layer height was kept constant across all grids to strictly enforce the $y^+ \approx 1$ requirement, while the bulk and inflation-layer resolutions were systematically refined. 

\begin{table}[htbp]
	\centering
	\footnotesize
	\caption{Wall-resolved channel mesh hierarchy used for the grid-independence study. The first-layer height is kept fixed to maintain $y^+\approx1$, while the streamwise resolution, wall-normal resolution, and inflation-layer resolution are refined.}
	\label{tab:appendix_channel_meshes}
	\renewcommand{\arraystretch}{1.2}
	\begin{tabularx}{\textwidth}{lccc}
		\hline
		\textbf{Parameter} & \textbf{Coarse} & \textbf{Medium} & \textbf{Fine} \\
		\hline
		Mesh label & \texttt{Coarse\_WallResolved} & \texttt{Medium\_WallResolved} & \texttt{Fine\_WallResolved} \\
		Streamwise cells & $10$ & $20$ & $40$ \\
		Wall-normal cells & $42$ & $51$ & $68$ \\
		First-layer height $\Delta y_1$ [m] & $1.709 \times 10^{-3}$ & $1.709 \times 10^{-3}$ & $1.709 \times 10^{-3}$ \\
		Minimum wall-normal spacing [m] & $1.709 \times 10^{-3}$ & $1.709 \times 10^{-3}$ & $1.709 \times 10^{-3}$ \\
		Maximum wall-normal spacing [m] & $7.707 \times 10^{-2}$ & $8.454 \times 10^{-2}$ & $7.589 \times 10^{-2}$ \\
		Inflation layers per wall & $10$ & $14$ & $18$ \\
		Growth ratio & $1.45$ & $1.35$ & $1.25$ \\
		Vertices & $473$ & $1{,}092$ & $2{,}829$ \\
		Triangular elements & $840$ & $2{,}040$ & $5{,}440$ \\
		\hline
	\end{tabularx}
\end{table}

Visual confirmation of this spatial convergence is provided by the wall-normal velocity profiles plotted in Figure~\ref{fig:channel_mesh_independence}. The \texttt{Fine\_WallResolved} grid was consequently selected for the final verification against the Ansys Fluent baseline, while the coarser grids demonstrate that the wall-normal velocity profile is already asymptotically converged.

\begin{figure}[htpb]
	\centering
	\includegraphics[width=0.6\textwidth]{example-image} % Replace with your channel mesh independence velocity profile plot
	\caption{Wall-normal velocity profiles for the Coarse, Medium, and Fine grids in the periodic channel flow benchmark. The overlapping profiles demonstrate asymptotic spatial convergence.}
	\label{fig:channel_mesh_independence}
\end{figure}

\section{Verification Benchmark 2: U-Bend Geometry}
\label{sec:app_ubend}

\subsection{Physical and Geometric Parameters}
\label{subsec:app_ubend_physical}
The second verification benchmark evaluates a highly convective U-bend. Because this FEniCS domain is a 2D planar approximation of the symmetry mid-plane of the 3D Ansys VMFL048 case, the 3D physical pipe diameter ($D$) is explicitly retained as the characteristic length scale. This ensures the Reynolds ($Re_D$) and Dean ($De$) numbers remain non-dimensionally consistent with the commercial baseline. The configuration is detailed in Table~\ref{tab:app_ubend_physical}.

\begin{table}[htbp]
	\centering
	\footnotesize
	\caption{Geometric, physical, and boundary-condition parameters for the U-bend benchmark.}
	\label{tab:app_ubend_physical}
	\renewcommand{\arraystretch}{1.2}
	\begin{tabularx}{\textwidth}{l X}
		\hline
		\textbf{Parameter} & \textbf{U-bend setting} \\
		\hline
		Benchmark type & Truncated 2D symmetry-plane U-bend based on Ansys VMFL048 \\
		Domain dimensions & Pipe diameter $D=0.028$ m; centreline bend radius $R_c=0.125$ m; straight-leg length $30D=0.840$ m \\
		Characteristic length & Physical pipe diameter $D=0.028$ m \\
		Reference / inlet velocity & Uniform prescribed inlet velocity $\mathbf{u}_{in}=(0,-1.42)$ m/s \\
		Kinematic viscosity & $\nu = 8.9 \times 10^{-7}$ m$^2$/s \\
		Reynolds number & $Re_D \approx 4.47 \times 10^4$ \\
		Dean number & $De \approx 1.50 \times 10^4$ \\
		Body force & $\mathbf{f}=(0,0)$ m/s$^2$ \\
		Driving condition & Prescribed inlet velocity and zero relative static pressure at the outlet \\
		Velocity boundary conditions & No-slip on all curved and straight pipe walls; natural velocity condition at pressure outlet \\
		Pressure boundary conditions & $p=0$ Pa at outlet; pressure left free at inlet and walls \\
		SA boundary conditions & $\tilde{\nu}_{in}=9.34 \times 10^{-5}$ m$^2$/s at inlet; $\tilde{\nu}=0$ at walls; natural condition at outlet \\
		Initial fields & $\mathbf{u}_0=(0,0)$ m/s, $p_0=0$ Pa, $\tilde{\nu}_0=\tilde{\nu}_{in}$ \\
		Verification mesh & \texttt{Medium\_WallResolved}; first layer $\Delta y_1 = 1.20 \times 10^{-5}$ m \\
		\hline
	\end{tabularx}
\end{table}

\subsection{Numerical Simulation Parameters}
\label{subsec:app_ubend_numerical}
Due to the adverse pressure gradients and secondary flow separation induced by the streamline curvature, the U-bend represents a significantly more demanding numerical challenge than the channel. Consequently, the FEniCS solver required a smaller pseudo-time step, tighter under-relaxation, and advanced iterative preconditioners to maintain stability, as detailed in Table~\ref{tab:app_ubend_numerical}.

\begin{table}[htbp]
	\centering
	\footnotesize
	\caption{Numerical simulation parameters for the highly convective FEniCS U-bend benchmark.}
	\label{tab:app_ubend_numerical}
	\renewcommand{\arraystretch}{1.2}
	\begin{tabularx}{\textwidth}{l c X}
		\hline
		\textbf{Parameter} & \textbf{Symbol} & \textbf{Setting} \\
		\hline
		Velocity space & $\mathbf{u}$ & Vector CG2 \\
		Pressure space & $p$ & Scalar CG1 \\
		SA working-variable space & $\tilde{\nu}$ & Scalar CG1 \\
		Quadrature degree & -- & $4$ \\
		Maximum outer Picard steps & $N_{\mathrm{Picard,max}}$ & $180$ \\
		SA solves per Picard step & -- & $1$ \\
		Outer velocity tolerance & $\epsilon_{\mathbf{u}}$ & $1.0 \times 10^{-4}$ \\
		Outer pressure tolerance & $\epsilon_p$ & $1.0 \times 10^{-4}$ \\
		Outer SA tolerance & $\epsilon_{\tilde{\nu}}$ & $1.0 \times 10^{-6}$ \\
		IPCS pseudo-time step & $\Delta t$ & $1.0 \times 10^{-4}$ s \\
		Maximum IPCS steps per Picard step & -- & $180$ \\
		IPCS velocity tolerance & -- & $1.0 \times 10^{-4}$ \\
		IPCS pressure tolerance & -- & $2.0 \times 10^{-3}$ \\
		Velocity relaxation factor & $\alpha_u$ & $0.3$ \\
		Pressure relaxation factor & $\alpha_p$ & $0.1$ \\
		SA relaxation factor & $\alpha_{\tilde{\nu}}$ & $0.01$ \\
		SA lower bound & $\tilde{\nu}_{\min}$ & $1.0 \times 10^{-12}$ \\
		Wall-distance equation & -- & Relaxed Eikonal, $\sigma_w=0.1$, $G_0=20$ \\
		Wall-distance Newton settings & -- & Relative tolerance $1.0 \times 10^{-8}$; absolute tolerance $1.0 \times 10^{-10}$; maximum $80$ iterations; relaxation $0.5$ \\
		Linear solver strategy & -- & IPCS blocks: GMRES with ILU velocity/correction preconditioning and Hypre-AMG pressure preconditioning; SA transport: default solver \\
		\hline
	\end{tabularx}
\end{table}

\subsection{Domain Truncation Study}
\label{subsec:app_ubend_truncation}
The original 3D Ansys VMFL048 benchmark utilizes highly extended straight legs measuring $1.555$ m to ensure fully developed flow. To reduce computational overhead for the topology optimization pipeline, these legs were truncated to $30D$ ($0.840$ m). A preliminary sensitivity study was conducted in Ansys Fluent, comparing the full-length and truncated geometries. To mathematically prove that this reduction does not artificially alter the flow physics, the cross-sectional velocity profiles at the critical bend exit were extracted for both domains, as plotted in Figure~\ref{fig:ubend_truncation_proof}. The complete overlap of the profiles confirms that the $30D$ truncation successfully isolated the curvature physics without artificially impacting the boundary layer separation.

\begin{figure}[htpb]
	\centering
	\includegraphics[width=0.6\textwidth]{example-image} % Replace with your U-bend truncation velocity comparison plot
	\caption{Velocity profile comparison at the U-bend exit between the full Ansys domain ($1.555$ m legs) and the truncated topology optimization domain ($30D$ legs). The complete overlap confirms the validity of the truncated geometry.}
	\label{fig:ubend_truncation_proof}
\end{figure}

\subsection{Spatial Verification and Grid Independence}
\label{subsec:app_ubend_mesh}
Managing the highly stretched, non-orthogonal inflation layers around the curved arcs of the U-bend required custom mesh generation via Gmsh. A spatial grid independence study was conducted across three refined unstructured grids. The node counts, element densities, and inflation layer parameters for these generated meshes are explicitly documented in Table~\ref{tab:appendix_ubend_meshes}. Similar to the channel benchmark, the wall-adjacent element height was strictly maintained at $1.20 \times 10^{-5}$ m to satisfy the low-Reynolds SA wall constraints, while the internal element sizing was systematically refined.

\begin{table}[htbp]
	\centering
	\footnotesize
	\caption{Wall-resolved U-bend mesh hierarchy used for the grid-independence study. The first-layer height is kept fixed to maintain $y^+\approx1$, while the bulk element size and inflation-layer resolution are refined.}
	\label{tab:appendix_ubend_meshes}
	\renewcommand{\arraystretch}{1.2}
	\begin{tabularx}{\textwidth}{lccc}
		\hline
		\textbf{Parameter} & \textbf{Coarse} & \textbf{Medium} & \textbf{Fine} \\
		\hline
		Mesh label & \texttt{Coarse\_WallResolved} & \texttt{Medium\_WallResolved} & \texttt{Fine\_WallResolved} \\
		Bulk element size [m] & $2.80 \times 10^{-3}$ & $1.40 \times 10^{-3}$ & $8.00 \times 10^{-4}$ \\
		First-layer height $\Delta y_1$ [m] & $1.20 \times 10^{-5}$ & $1.20 \times 10^{-5}$ & $1.20 \times 10^{-5}$ \\
		Inflation layers & $10$ & $14$ & $18$ \\
		Growth ratio & $1.45$ & $1.30$ & $1.25$ \\
		Inflation thickness [m] & $1.07 \times 10^{-3}$ & $1.53 \times 10^{-3}$ & $2.62 \times 10^{-3}$ \\
		Vertices & $23{,}751$ & $73{,}978$ & $182{,}450$ \\
		Triangular elements & $45{,}998$ & $144{,}952$ & $359{,}646$ \\
		\hline
	\end{tabularx}
\end{table}

Convergence was visually verified by extracting the velocity profiles directly at the bend exit, a highly sensitive region for curvature-induced separation, as illustrated in Figure~\ref{fig:ubend_mesh_independence}. The chosen \texttt{Medium\_WallResolved} grid perfectly captured the velocity gradients near the wall and provided sufficient isotropic resolution in the bulk flow to resolve the separation bubbles without numerical instability.

\begin{figure}[htpb]
	\centering
	\includegraphics[width=0.6\textwidth]{example-image} % Replace with your U-bend mesh independence velocity profile plot
	\caption{Velocity profiles at the U-bend exit for the Coarse, Medium, and Fine grids. The profiles demonstrate spatial convergence, justifying the selection of the Medium grid for the final verification.}
	\label{fig:ubend_mesh_independence}
\end{figure}
